\documentclass[a4paper,12pt]{article}
\usepackage{graphicx}
\usepackage[utf8]{inputenc}  
\usepackage[T1]{fontenc}
\usepackage{gensymb}

\usepackage{tgheros,textcomp}%Police Arial
\renewcommand{\familydefault}{\sfdefault}

\usepackage[top = 1cm, bottom = 2cm, left = 1cm, right = 1cm]{geometry}
\usepackage{titlesec}
\setcounter{secnumdepth}{4}
\usepackage{xcolor}
\usepackage{amssymb} %double flèche de réaction
\usepackage{xtab}%tableau sur plusieurs pages
\usepackage{esvect}%grands vecteurs

\usepackage[version=4]{mhchem} %equations chimiques

\usepackage{array}
\usepackage{chemfig}
\usepackage{multirow}
\usepackage{sectsty}
\usepackage{caption}
\usepackage{amsmath}
\usepackage{amssymb}
\usepackage[cal=boondoxo]{mathalfa} 
\usepackage{enumitem}
\usepackage{lastpage}
\usepackage{tcolorbox}
\usepackage{siunitx}
\usepackage{tikz}
\usetikzlibrary{calc}
\usepackage{multicol}
\usepackage{eqnarray}

\sisetup{
	detect-all,
	output-decimal-marker={,},
	group-minimum-digits = 3,
	group-separator={~},
	number-unit-separator={~},
	inter-unit-product={~}
} %setup pour SIunitx

\usepackage{listings}%pour insérer le code python
\definecolor{codegreen}{rgb}{0,0.6,0}
\definecolor{codegray}{rgb}{0.5,0.5,0.5}
\definecolor{codepurple}{rgb}{0.58,0,0.82}
\definecolor{backcolour}{rgb}{0.95,0.95,0.92}

\sectionfont{\fontsize{14}{8}\selectfont}
\renewcommand{\thesection}{\arabic{section}.}
\renewcommand{\thesubsection}{\hspace{1cm}\alph{subsection}.}

\title{\vspace{-1cm}\large 
	\begin{tabular}{|m{7cm}|>{\centering}m{6cm}|>{\centering\arraybackslash}m{4cm}|}
		\hline
		 Nom :& \textbf{\textsc{2GT DS N°2}} &  Note\\
		 &&\\
		 Prénom : &&\\
		\hline
\end{tabular}}
\date{}

\newpagestyle{mystyle}{%
	\setfoot{}{\thepage\,$|$\,\pageref{LastPage}}{}
}
\renewpagestyle{plain}{%
	\setfoot{}{\thepage\,$|$\,\pageref{LastPage}}{}
}

\pagestyle{mystyle}

\newenvironment{itemize*}%
{\begin{itemize}[label=$-$]%
		\setlength{\itemsep}{0pt}%
		\setlength{\parskip}{0pt}}%
	{\end{itemize}}

\newenvironment{itemize**}%
{\begin{itemize}[label=]%
		\setlength{\itemsep}{0pt}%
		\setlength{\parskip}{0pt}}%
	{\end{itemize}}

\begin{document}
	\maketitle
	\vspace{-3cm}
\section*{Exercice 1 : Sons}
\begin{tabular}{cc}
	\includegraphics[scale=.5]{Ressources/2GT_C7_Eval_Fig1.png}&\includegraphics[scale=.5]{Ressources/2GT_C7_Eval_Fig3.png}\\
	\includegraphics[scale=.5]{Ressources/2GT_C7_Eval_Fig4.png}&\includegraphics[scale=.5]{Ressources/2GT_C7_Eval_Fig2.png}\\
\end{tabular}
\begin{enumerate}
	\item Dire si les signaux sonores présentés ci dessus sont périodiques. Justifier chaque réponse. (1 pt)
	\item Pour un signal périodique, définir un motif. Repasser en couleur un motif pour chaque signal périodique trouvé dans la réponse précédente. (2 pts)
	\item Déterminer les fréquences des signaux précédent si possible (des calculs détaillés sont attendus à chaque fois). (3 pts)
\end{enumerate}

\section*{Exercice 2 : Écholocalisation}
L'écholocalisation est un système de perception de l'environnement développé par plusieurs espèces (chauve-souris, cétacés) afin de percevoir l'environnement proche dans utiliser la vision.\\
Il s'agit pour l'animal de produire un son de haute fréquence qui va se propager jusqu'à un obstacle puis revenir à l'oreille de l'animal. L'animal joue ainsi, à la fois le rôle d'émetteur et de récepteur. La perception de proximité est faite par la mesure du temps mis par l'onde sonore pour revenir à son émetteur.
\begin{enumerate}
	\item Proposer un schéma permettant de représenter le système d'écholocalisation. Faire figurer un émetteur ($E$), un objet à localiser ($A$), un récepteur ($R$) et la distance entre le système émetteur/récepteur et l'objet ($d$). (1 pt)
	\item Dans le cas d'une chauve-souris, exprimer le temps ($\Delta t_{CS}$) mis par l'onde sonore pour revenir vers elle après avoir rencontré un insecte situé à 3 m. (1pt)
	\item Sachant que, le son a une célérité $c_{air}$ = \SI{340}{\meter\per\second}, calculer la valeur de $\Delta t_{CS}$. (1 pt)
	\item Dans le cas où l'animal est un dauphin, exprimer le temps ($\Delta t_{D}$) mis par l'onde sonore pour revenir vers elle après avoir rencontré un saumon situé à 3 m. Puis calculer $\Delta t_{D}$, sachant que $c_{eau}$ = \SI{1500}{\meter\per\second}. (2 pts)
	\item Lequel des deux animaux a le potentiel d'être le plus rapide pour chasser ? (1 pt)
	%\item \{Question supplémentaire non notée :\} imaginer d'autre paramètres pouvant influencer la vitesse en chasse des deux animaux
\end{enumerate}

\section*{Exercice 3 : Atomes}
On considère deux noyaux atomiques : le bismuth 209 (\ce{^209_82 Bi}) et le polonium 209 (\ce{^209_84 Po}).
\begin{enumerate}
	\item Donner la composition des deux noyaux et les comparer. (2 pts)
	\item Calculer la masse des deux noyaux et les comparer.  (2 pts)
	\item Donner dans chacun des cas le nombre d'électrons présent pour les atomes associés.  (1 pt)
	\item S'agit-il d'isotopes ? Justifier.  (1 pt)
\end{enumerate}

\textbf{Données :}
\begin{itemize*}
	\item $m_{\text{nucléon}}$ = \SI{1.67E-27}{kg}
	\item $m_{\text{électron}}$ = \SI{9.11E-31}{kg}
\end{itemize*}
\end{document}